\exercisesection{Spectrum of Endomorphisms of Banach Spaces}

In the exercises below, all Banach spaces are over C.

\begin{enumerate}
\def\labelenumi{\arabic{enumi})}
\tightlist
\item
  Let \(\mathrm{E} = \mathcal{C}([0,1])\) and let \(u\) be the endomorphism of \(\mathrm{E}\) defined by
\end{enumerate}

\[
u (f) (t) = \int_ {0} ^ {t} f (t) d t.
\]

\begin{enumerate}
\def\labelenumi{\alph{enumi})}
\item
  The spectral radius of \(u\) is zero. (Estimate the norm of \(u^n\) .)
\item
  The spectrum of \(u\) consists only of 0, and, for every \(n \geqslant 1\), the kernel of \(u^{n}\) consists only of 0.
\end{enumerate}

(In particular, 0 is an isolated point of the spectrum of \(u\), but the closure of the union of the kernels of \(u^n\), for \(n \geqslant 1\), is not equal to the spectral subspace \(\mathrm{E}_0(x)\).)

\begin{enumerate}
\def\labelenumi{\arabic{enumi})}
\setcounter{enumi}{1}
\item
  There exists a Banach space E and an endomorphism u of E such that 0 is an eigenvalue of u but 0 is not isolated in \(\mathrm{Sp}(u)\).
\item
  Give an example of a Hilbert space E and an endomorphism \(u \in \mathcal{L}(\mathrm{E})\) such that 0 is an eigenvalue of u but not of \(u^{*}\).
\item
  Let \(E = \ell^{2}(\mathbf{Z})\) and let \(u\) be the endomorphism such that \(u((\lambda_{n})) = (\lambda_{n+1})\).
\end{enumerate}

\begin{enumerate}
\def\labelenumi{\alph{enumi})}
\item
  Show that \(u\) is unitary.
\item
  Compute the spectrum of \(u\).
\item
  The finite spectrum of \(u\) is empty.
\end{enumerate}

\begin{enumerate}
\def\labelenumi{\arabic{enumi})}
\setcounter{enumi}{4}
\tightlist
\item
  Let \(s\) be the endomorphism of \(\ell^2 (\mathbf{N})\) such that
\end{enumerate}

\[
s (a _ {0}, a _ {1}, \dots , a _ {n}, \dots) = (a _ {1}, \dots , a _ {n}, \dots).
\]

Show that the numerical range of \(s\) is the open unit disk \(\Delta\) in C, and that the spectrum of \(s\) is the closed unit disk in C. Determine the eigenvalues of \(s\). Do the same for the adjoint \(s^{*}\) of \(s\).

\begin{enumerate}
\def\labelenumi{\arabic{enumi})}
\setcounter{enumi}{5}
\item
  Let \(u \in \mathcal{L}(\mathbf{C}^n)\) be a trace-zero endomorphism. There exists an orthonormal basis B of \(\mathbf{C}^n\) such that all diagonal entries of the matrix representing \(u\) in the basis B are zero. (Prove that 0 belongs to the numerical range of \(u\).)
\item
  Let \(u\) be an endomorphism of a complex Hilbert space E such that the spectrum of \(u\) is contained in the open unit disk in C. Let F be the Hilbert space \(\ell_{\mathrm{E}}^{2}(\mathbf{N})\) of square-summable sequences of elements of E (EVT, V, p.~18, example), and let \(s_{\mathrm{F}}\) be the endomorphism of F such that
\end{enumerate}

\[
s _ {\mathrm{F}} (a _ {0}, a _ {1}, \ldots , a _ {n}, \ldots) = (a _ {1}, \ldots , a _ {n} \ldots).
\]

There exists a closed subspace H of F and a continuous bijective linear map \(v: E \rightarrow H\) such that \(u = v^{-1}s_{F}v\).

\begin{enumerate}
\def\labelenumi{\arabic{enumi})}
\setcounter{enumi}{7}
\tightlist
\item
  Let \(u\) be an endomorphism of a complex Hilbert space E. Let \(S_u\) be the set of endomorphisms of E of the form \(wuw^{-1}\), where \(w\) ranges over the set of invertible endomorphisms of E.
\end{enumerate}

\begin{enumerate}
\def\labelenumi{\alph{enumi})}
\tightlist
\item
  We have
\end{enumerate}

\[
\varrho (u) = \inf _ {v \in \mathrm{S} _ {u}} \| v \|.
\]

(One may assume that \(\varrho(u) \leqslant 1\); apply the preceding exercise to the endomorphisms \(u_{\varepsilon} = (1 + \varepsilon)^{-1} u\) for \(\varepsilon > 0\).)

\begin{enumerate}
\def\labelenumi{\alph{enumi})}
\setcounter{enumi}{1}
\tightlist
\item
  The convex hull of the spectrum of \(u\) is equal to the intersection of the subsets \(\iota(v)\), for \(v \in S_u\) ("Hildebrandt's theorem").
\end{enumerate}

¶ 9) Let E be a complex Banach space.

\begin{enumerate}
\def\labelenumi{\alph{enumi})}
\item
  Let \(u \in \mathcal{L}(\mathrm{E})\). Consider the endomorphisms \(\gamma_{u}\) and \(\delta_{u}\) of \(\mathcal{L}(\mathrm{E})\) defined by \(\gamma_{u}(v) = u \circ v\) and \(\delta_{u}(v) = v \circ u\). We have \(\operatorname{Sp}(u) = \operatorname{Sp}(\gamma_{u}) = \operatorname{Sp}(\delta_{u})\), where the last two spectra are relative to the algebra \(\mathcal{L}(\mathcal{L}(\mathrm{E}))\).
\item
  For every function \(f\) holomorphic in a neighborhood of the spectrum of \(u\), one has \(\boldsymbol{\gamma}_{f(u)} = f(\boldsymbol{\gamma}_{u})\) and \(\boldsymbol{\delta}_{f(u)} = f(\boldsymbol{\delta}_{u})\).
\item
  Let \(u_{1}\) and \(u_{2}\) be in \(\mathcal{L}(\mathrm{E})\). Let \(\varrho\) be the endomorphism \(v \mapsto u_{1} \circ v \circ u_{2}\) of \(\mathcal{L}(\mathrm{E})\). Its spectrum is equal to \(\operatorname{Sp}(u_{1})\operatorname{Sp}(u_{2}) \subset \mathbf{C}\).
\end{enumerate}

\(d)\) Let \(u_{1}\) and \(u_{2}\) be in \(\mathcal{L}(\mathrm{E})\). Let \(\varrho\) be the endomorphism \(v \mapsto u_{1} \circ v \circ u_{2}\) of the space \(\mathcal{L}_{2}(\mathrm{E})\) of Hilbert--Schmidt maps from E to E. Its spectrum is equal to \(\mathrm{Sp}(u_{1})\mathrm{Sp}(u_{2}) \subset \mathbf{C}\). (Apply a similar method.)

\begin{enumerate}
\def\labelenumi{\arabic{enumi})}
\setcounter{enumi}{9}
\item
  Let E be a complex Hilbert space. Let \(x_1\) and \(x_2\) be endomorphisms of E. Show that the spectrum of the endomorphism \(x_1 \otimes x_2\) of E \(\widehat{\otimes}_2\) E is equal to \(\mathrm{Sp}(x_1) \mathrm{Sp}(x_2)\). (Identify \(\overline{\mathrm{E}} \widehat{\otimes}_2\) E with the space of Hilbert--Schmidt operators from E to E, cf.~EVT 5, p.~52, th. 1, b), and apply the preceding exercise.)
\item
  Let \(u\) be an endomorphism of a complex Hilbert space \(E\) and let \((j, |u|)\) be its polar decomposition.
\end{enumerate}

\begin{enumerate}
\def\labelenumi{\alph{enumi})}
\item
  For \(j\) to commute with \(|u|\), it is necessary and sufficient that \(u\) commute with \(u^{*}u\).
\item
  Prove that one then has \(\|u^{*}(x)\|\leqslant\|u(x)\|\) for all \(x\in E\), and that \(\varrho(u)=\|u\|\).
\item
  Let \(E = \ell^{2}(\mathbf{N})\) and let \(u\) be the endomorphism of E given by
\end{enumerate}

\[
u ((a _ {0}, a _ {1}, \dots)) = (0, a _ {0}, a _ {1}, \dots).
\]

It is not normal. Compute the polar decomposition \((j, |u|)\) of u; \(j\) commutes with \(|u|\).

\begin{enumerate}
\def\labelenumi{\arabic{enumi})}
\setcounter{enumi}{11}
\tightlist
\item
  Let \(u\) be an endomorphism of a Banach space E. We denote by \(\sigma_{a}(u)\) the set of \(\lambda\in\mathbf{C}\) such that, for every \(\alpha\in\mathbf{R}_{+}^{*}\), there exists \(x\in\mathrm{E}\) such that \(\|\lambda x-u(x)\|<\alpha\|x\|\).
\end{enumerate}

\begin{enumerate}
\def\labelenumi{\alph{enumi})}
\item
  Prove that \(\sigma_{a}(u)\) is a closed subset of \(\sigma(u)\).
\item
  Prove that \(\sigma_{a}(u)\) contains the boundary of \(\sigma(u)\).
\end{enumerate}

\begin{enumerate}
\def\labelenumi{\arabic{enumi})}
\setcounter{enumi}{12}
\tightlist
\item
  Let E be a complex Hilbert space. Let N be a map from \(\mathcal{L}(\mathrm{E})\) into \(\mathfrak{P}(\mathbf{C})\) such that
\end{enumerate}

\begin{enumerate}
\def\labelenumi{\alph{enumi})}
\item
  \(N(u)\) is compact for every \(u \in \mathcal{L}(\mathrm{E})\);
\item
  \(\mathrm{N}(\alpha u + \beta 1_{\mathrm{E}}) = \alpha \mathrm{N}(u) + \beta\) for all \(u \in \mathcal{L}(\mathrm{E})\) and every \((\alpha, \beta) \in \mathbf{C}^{2}\);
\item
  \(\mathrm{N}(u)\subset\{z\in\mathbf{C}\mid\mathcal{R}(z)\geqslant0\}\) if and only if \(u+u^{*}\) is positive.
\end{enumerate}

  Then \(\mathrm{N}(u)=\overline{\iota(u)}\) for all \(u\in\mathcal{L}(\mathrm{E})\). (Show that \(u\mapsto\overline{\iota(u)}\) has the indicated properties; let \(N_{1}\) and \(N_{2}\) be maps satisfying them; assuming that \(\mathrm{N}_{1}(u)\) is not contained in \(\mathrm{N}_{2}(u)\), find \(\alpha\) and \(\beta\) in C such that \(\mathrm{N}_{2}(\alpha u+\beta)\) is contained in the half-plane \(\mathcal{R}(z)\geqslant0\) but \(\mathrm{N}_{1}(\alpha u+\beta)\) is not, and conclude a contradiction using c).)

\begin{enumerate}
\def\labelenumi{\arabic{enumi})}
\setcounter{enumi}{13}
\item
  Give an example of a complex Hilbert space and of \(u \in \mathcal{L}(\mathrm{E})\) such that \(\iota(u)\) is not closed.
\item
  The set X of nonempty compact subsets of C is equipped with the distance defined in TG, IX, p.~91, exerc. 6. Let E be a complex Hilbert space.
\end{enumerate}

\begin{enumerate}
\def\labelenumi{\alph{enumi})}
\item
  The map \(u \mapsto \overline{\iota(u)}\) from \(\mathcal{L}(\mathrm{E})\) to X is continuous. If E is infinite-dimensional, it is not continuous when \(\mathcal{L}(\mathrm{E})\) is equipped with the topology of simple convergence.
\item
  The map \(u \mapsto \operatorname{Sp}(u)\) from \(\mathcal{L}(\mathrm{E})\) to X is not continuous.
\end{enumerate}

\begin{enumerate}
\def\labelenumi{\arabic{enumi})}
\setcounter{enumi}{15}
\tightlist
\item
  Let X be a locally compact space, \(\mu\) a positive measure of total mass 1 on X, and \(f: X \to X\) a map such that \(f(\mu) = \mu\).
\end{enumerate}

\begin{enumerate}
\def\labelenumi{\alph{enumi})}
\item
  The map \(\varphi \mapsto \varphi \circ f\) induces a unitary endomorphism \(u_{f}\) of \(\mathrm{L}^2 (\mathrm{X},\mu)\).
\item
  The real number 1 is an eigenvalue of \(u_f\) of multiplicity 1 if and only if, for every measurable subset A of X such that \(f(A) = A\), one has \(\mu(A) \in \{0,1\}\). One then says that \(f\) is \(\mu\)-ergodic.
\item
  Let \(\mathrm{E}\subset\mathrm{L}^{2}(\mathrm{X},\mu)\) be the eigenspace of \(u_{f}\) for the eigenvalue 1. The sequence \((v_{\mathrm{N}})_{\mathrm{N}\in\mathbf{N}}\) defined by
\end{enumerate}

\[
v _ {\mathrm{N}} = \frac {1}{\mathrm{N}} \sum_ {n = 0} ^ {\mathrm{N-1}} u _ {f} ^ {n}
\]

converges to the orthogonal projection onto E in the space \(\mathcal{L}(\mathrm{L}^{2}(\mathrm{X},\mu))\) equipped with the topology of simple convergence ("von Neumann's ergodic theorem"). (Note that the orthogonal complement of E is the closure of the image of \(u_{f}-1_{\mathrm{E}}\); then verify that \(v_{\mathrm{N}}(\varphi)\to0\) when \(\varphi\in\mathrm{E}^{\circ}\).)

\(d)\) One defines an endomorphism \(\widetilde{u}_f\) of \(\mathrm{L}^1 (\mathrm{X},\mu)\) as in part \(a)\), and we set

\[
\widetilde {v} _ {\mathrm{N}} = \frac {1}{\mathrm{N}} \sum_ {n = 0} ^ {\mathrm{N-1}} \widetilde {u} _ {f} ^ {n}.
\]

For every \(\varphi\in\mathrm{L}^{1}(\mathrm{X},\mu)\) the sequence \((\widetilde{v}_{\mathrm{N}}(\varphi))_{\mathrm{N}\in\mathbf{N}}\) converges in \(\mathrm{L}^{1}(\mathrm{X},\mu)\) to an element \(\widetilde{\varphi}\) of the eigenspace of \(\widetilde{u}_{f}\) for the eigenvalue 1. (First consider the case where \(\varphi\) is bounded, and apply the previous question.)

\begin{enumerate}
\def\labelenumi{\alph{enumi})}
\setcounter{enumi}{4}
\tightlist
\item
  Let I be a finite set and X = I \(^{Z}\) endowed with the product measure of the measures \(\mu_n\) on I such that \(\mu_n(i) = 1/\text{Card(I)}\) for all \(i \in \text{I}\). Set \(f((i_n)) = (i_{n+1})\) for all \((i_n) \in \text{X}\). Then \(f(\mu) = \mu\) and \(f\) is \(\mu\)-ergodic (“Bernoulli shift”).
\end{enumerate}

\(f)\) Let \(\mathrm{X} = \mathbf{R} / \mathbf{Z}\) and \(\mu\) the normalized Haar measure on \(\mathrm{X}\). Let \(a \in \mathrm{X}\) and set \(f(x) = x + a\). Then \(f(\mu) = \mu\) and \(f\) is \(\mu\)-ergodic if and only if \(a\) is the class of an irrational number.

